\chapter{Radar Physics and Thematic Literature Review}
\label{ch:physics_related_work}

This chapter develops the signal model and methodological background for the thesis. It connects
the physics of continuous-wave (CW) quadrature radar to the processing choices made in the later
experimental chapters: complex demodulation, time--frequency artifact separation, harmonic rate
estimation, temporal decoding, selective prediction, and clustered evaluation. The discussion is
organized around six research themes represented in the repository's curated bibliography and
literature matrix, which together contain 71 references. The matrix records each entry's method,
validation context, and relevance to the work packages. The repository does not include a
reproducible database-search log, screening counts, or duplicate-review record for a new systematic
review. Accordingly, this chapter is a structured thematic synthesis of that fixed source set; it
does not claim to report a newly executed PRISMA-style systematic search. The 60-study systematic
review by Senechal et al. is discussed as a published review of the field, not as a search performed
for this thesis \citep{senechal2026next}.

The chapter also establishes a boundary between physical quantities, measured signals, and
algorithmic estimates. The displacement of the chest wall is denoted by $\chestDisp$; the
mean-centered complex baseband is $\csdSignal$; and the filtered signal supplied to the rate
estimator is $\chestDispEst$. These objects have different meanings and units. A voltage-like
complex radar return is not a calibrated displacement measurement, and a rate estimate is not a
clinical alarm. The derivations below explain what the single-channel sensor can encode and where
its information is intrinsically incomplete.

\section{Continuous-Wave Quadrature Radar Physics}
\label{sec:ch2_radar_physics}

\subsection{Propagation, round-trip phase, and the quadrature observation}

A monostatic CW radar transmits a sinusoid at carrier frequency $\carrierFreq$ and observes the
complex echo after reflection from the infant's thorax. For a dominant reflector at nominal
standoff $\nomRange$ undergoing a small radial displacement $\chestDisp$, the propagation path
changes by approximately twice the displacement: once on transmission and once on return. The
corresponding phase can be written
\begin{equation}
  \phaseSig = \phaseZero + \frac{4\pi}{\radarLambda}\chestDisp,
  \qquad \radarLambda = \frac{\speedLight}{\carrierFreq},
  \label{eq:ch2_radar_phase}
\end{equation}
where $\phaseZero$ includes the nominal range phase and fixed radio-frequency phase offsets. At
$\carrierFreq=24\,\mathrm{GHz}$, the free-space wavelength is approximately
$\radarLambda=12.5\,\mathrm{mm}$. The phase sensitivity is therefore
$4\pi/\radarLambda$, or about $1.01$ radians per millimeter of radial motion. A one-millimeter
displacement therefore produces an approximately one-radian phase change in this idealized
geometry. This sensitivity explains why small physiological movements can modulate a coherent
return, including through clothing or non-metallic incubator surfaces; it does not establish that
the radar can detect every movement under arbitrary placement, attenuation, or clutter
\citep{droitcour2004range,li2013review,hu2014noncontact}.

After downconversion and low-pass filtering, an idealized quadrature receiver observes
\begin{align}
  \iqI &= A_I\cos(\phaseSig)+\iqDcI+\nu_I(t),\nonumber\\
  \iqQ &= A_Q\sin(\phaseSig+\iqPhaseSkew)+\iqDcQ+\nu_Q(t),
  \label{eq:ch2_iq_model}
\end{align}
where $A_I$ and $A_Q$ are effective channel amplitudes, including return amplitude and receiver
gain; $\iqDcI$ and $\iqDcQ$ are receiver offsets,
$\iqPhaseSkew$ is departure from ideal quadrature, and $\nu_I,\nu_Q$ collect receiver noise
and unmodeled interference. The complex observation is $\iqComplex=\iqI+j\iqQ$. In the ideal
case, its samples trace a circle as phase varies. Gain imbalance, quadrature skew, and DC offsets
turn that circle into a shifted ellipse. This geometric interpretation is useful because it
separates three phenomena that are often conflated: physiological phase modulation, receiver
calibration error, and additive interference. Ellipse fitting can estimate a receiver transform,
but short arcs and weak motion make that estimate poorly conditioned; more elaborate calibration
is not automatically more reliable in the low-motion windows that matter most
\citep{park2007arctangent,singh2013data,zakrzewski2014quadrature,fitzgibbon1999direct}.

The single-channel CW architecture has an important information limit. It observes the sum of
reflections in its field of view but does not encode an independent time-of-flight measurement.
It therefore cannot, by itself, separate two reflectors at different ranges. Static reflections
also contribute to the complex sum and set the operating phase $\phaseZero$. A sufficiently large
change in range, a deep fade, receiver saturation, or another moving reflector can distort the
baseband trajectory. The simple phase model is a useful local approximation for small motion; it
does not imply that every periodic component in the measured signal originates in respiration.

\subsection{Jacobi--Anger expansion and harmonic structure}

Respiratory displacement is not exactly sinusoidal, but a sinusoid is a tractable starting point
for understanding the radar observation. Let
\begin{equation}
  \chestDisp = A_r\sin(2\pi f_b t),
  \qquad \modIndex = \frac{4\pi A_r}{\radarLambda},
  \label{eq:ch2_modulation_index}
\end{equation}
where $A_r$ is displacement amplitude and $f_b$ is the respiratory fundamental frequency. After
removing a constant complex gain, the ideal phase-modulated return is proportional to
$\exp[j\modIndex\sin(2\pi f_b t)]$. The Jacobi--Anger identity gives
\begin{equation}
  \exp[j\modIndex\sin(2\pi f_b t)]
  = \sum_{n=-\infty}^{\infty}
      \besselJ{n}{\modIndex}\exp(j2\pi n f_b t).
  \label{eq:ch2_jacobi_anger}
\end{equation}
The return thus contains components at integer multiples of the breathing frequency, weighted by
Bessel functions. These components are a consequence of phase modulation; their existence does
not require random body movement or a defective sensor. Their relative magnitudes depend on
$\modIndex$, so the fundamental need not be the largest spectral peak. Radar vital-sign work has
long treated this nonlinear mapping and its harmonic consequences as central to demodulation and
rate estimation \citep{li2008random,morgan2009novel,wu2013phase,li2013review,tu2013respiration}.

If chest displacement contains both respiratory and cardiac terms, the phase becomes a sum of
modulations. The complex exponential then factors into a respiratory series and a cardiac series;
their product contains combination frequencies $n f_b + m f_h$, where $f_h$ is a cardiac
frequency and $n,m$ are integers. Realistic displacement also departs from a pure sinusoid, adding
further harmonic energy. A component at twice the fundamental is predicted by the ideal phase
modulation model, whereas a subharmonic is not: a subharmonic rate estimate can instead arise from
period ambiguity, alternating or period-doubled cycles, or other signal components. A nearby
cardiac--respiratory combination can likewise be spectral energy without being a separate
respiratory rate. Harmonic cancellation and cardiopulmonary state-space approaches address related
ambiguity, but their assumptions differ in how they represent the source waveform and motion
clutter \citep{naishadham2016estimation,beltrao2023adaptive}.

\subsection{Complex signal demodulation and calibrated alternatives}

An arctangent demodulator attempts to recover phase directly from the quadrature angle. In an ideal
receiver with its center removed, this can be expressed schematically as
$\hat{\phi}_{\mathrm{AD}}(t)=\operatorname{unwrap}(\operatorname{atan2}(Q(t),I(t)))$; multiplying
by $\radarLambda/(4\pi)$ yields a displacement-like phase trace. The operation is vulnerable
when the signal passes near the origin, where small noise changes can cause large phase changes,
and when offsets or ellipse distortion are not corrected. Differentiation-and-cross-multiplication
(DACM) instead estimates phase increments from adjacent quadrature samples and integrates them,
avoiding an explicit phase unwrap, but it too depends on suitable center and imbalance handling.
Methods for DC compensation, ellipse fitting, Gram--Schmidt correction, and extended DACM address
these receiver effects in different ways \citep{park2007arctangent,wang2014noncontact,fitzgibbon1999direct,singh2013data,zakrzewski2014quadrature}.

The locked baseline used in this thesis retains both channels and subtracts their per-window means:
\begin{equation}
  \csdSignal = (b_I-\bar b_I)+j(b_Q-\bar b_Q).
  \label{eq:ch2_csd}
\end{equation}
This complex signal demodulation (CSD) avoids converting the trajectory to an angle before
downstream processing. Mean subtraction removes a constant offset over the analysis window, while
the complex representation preserves the orientation and relative movement of the I/Q trajectory.
If a centered signal is transformed by a fixed complex gain and rotation,
$x'(t)=a\exp(j\vartheta)x(t)$ for $a>0$, then mean-centering commutes with that transform:
$x'_c(t)=a\exp(j\vartheta)x_c(t)$. This is an equivariance property of the representation, not
invariance of every possible downstream statistic: amplitudes still scale by $a$, so a later
algorithm must normalize scale or use scale-invariant features when that is required. CSD also does
not correct time-varying offsets, severe ellipse distortion, clipping, or multiple-reflector
interference. These limits motivate evaluating calibration variants separately rather than
assuming that a more complex demodulator must improve rate estimates
\citep{li2008random,li2013review,droitcour2004range}.

\section{Contactless Neonatal Monitoring and Radar Modalities}
\label{sec:ch2_nicu_literature}

The clinical literature spans several sensor families that trade spatial information, hardware
complexity, and dependence on the environment. Single-channel CW radar is attractive for its
compact front end and sensitivity to small periodic displacement, and it can operate without
attaching an electrode or optical marker to the infant. The 24 GHz NICU study by Beltrão et al.
established a prominent benchmark pipeline combining quadrature processing, short-time Fourier
analysis, non-negative matrix factorization (NMF) for random-body-movement mitigation, filtering,
and harmonic nonlinear least squares (NLS). Its results demonstrated feasibility under a range of
recording conditions, while also documenting interference and reference-device saturation as
practical issues. The paper reports 18 retained measurements and approximately 4,964 valid rate
estimates after its exclusions. Those published figures provide historical context; they do not
resolve the local cohort identifiers, exact window denominator, reference construction, or solver
details required for an exact reproduction claim in the present repository
\citep{beltrao2022contactless}.

Other work illustrates that “contactless radar” is not a single measurement design. Impulse-radio
ultra-wideband (IR-UWB) systems use short broadband pulses and can exploit range profiles to
separate chest motion from other reflections. Neonatal studies have compared IR-UWB-derived
respiration with clinical references and have examined movement censoring and feasibility in the
NICU \citep{kim2019noncontact,lee2020feasibility}. Frequency-modulated continuous-wave (FMCW)
systems use a swept carrier to encode range in beat frequency; multiple antennas or spatial
processing can add information that a single CW channel lacks. Recent work has compared mmWave
FMCW and RGB sensing in preterm infants, while a separate low-cost FMCW protocol describes a
clinical monitoring study design. Controlled comparisons of IR-UWB and FMCW further show why
range resolution, clutter, signal-to-noise ratio, and the task definition should be considered
together \citep{liao2025noncontact,hari2025contactless,wang2020experimental}. A synchronous
time-of-flight camera and microwave radar study represents a multimodal strategy: optical depth
information can help define a region of interest, while radar measures phase-sensitive motion
\citep{gleichauf2021automated}.

{\sloppy
Separate work on joint waveform and receiver design addresses detection under signal-dependent
interference, illustrating that some limitations arise at the sensing front end as well as in
post-processing \citep{beltrao2020joint}. These architectures do not form a universal ranking.
Range-resolving and multimodal devices may
provide additional ways to localize motion, but they add hardware, calibration, and processing
requirements. A single CW channel is simpler, yet its output is more vulnerable to superposed
reflections and cannot itself determine which range produced a component. Clinical performance
also depends on placement, infant posture, clothing, incubator geometry, motion, and the reference
used to define respiratory rate. Adult radar datasets and general engineering demonstrations can
inform signal processing but cannot substitute for neonatal validation; the synchronized radar
dataset of Schellenberger et al. is an adult clinical resource, whereas the Gleichauf et al.
thorax-simulator dataset is an external physical testbed rather than a cohort of monitored infants
\citep{schellenberger2020dataset,gleichauf2024time}.
\par}

At the field level, the systematic review by Senechal et al. included 60 observational studies of
non-contact and wireless vital-sign monitoring in neonatal intensive care. It reported a median
sample size of 10 patients and emphasized the need for longer recordings and more transparent
reporting of study design and participant selection. A broader clinical-technology review by
Krbec et al. likewise places non-contact monitoring within the unresolved problem of reliable
clinical surveillance rather than treating sensor feasibility as evidence of improved outcomes
\citep{senechal2026next,krbec2024emerging}. Together, this literature motivates three separate
questions: whether a sensor can detect respiratory movement, whether an algorithm agrees with a
chosen reference across representative conditions, and whether its outputs improve a clinical
decision. Evidence for the first question does not by itself answer the second or third.

\section{Spectrograms and Artifact-Separation Methods}
\label{sec:ch2_separation}

\subsection{Short-time Fourier representation}

The short-time Fourier transform (STFT) represents local spectral content over overlapping frames.
For a discrete signal $x[n]$, analysis window $w[n]$, frame hop $H_{\mathrm{stft}}$, and transform
length $N_{\mathrm{FFT}}$, one convention is
\begin{equation}
  X[k,m]=\sum_n x[n]w[n-mH_{\mathrm{stft}}]
  \exp\left(-j\frac{2\pi kn}{N_{\mathrm{FFT}}}\right).
  \label{eq:ch2_stft}
\end{equation}
The time--frequency trade-off is direct: longer windows sharpen frequency discrimination but blur
rapid changes; shorter windows localize transients but broaden spectral peaks. The locked v1
implementation uses a rectangular STFT with $N_{\mathrm{win}}=48$ samples, hop
$H_{\mathrm{stft}}=16$ samples, and $N_{\mathrm{FFT}}=256$ at a 16 Hz slow-time sampling rate.
These values correspond to a 3 s frame, a 1 s frame advance, and a sampled FFT-bin spacing of
0.0625 Hz (3.75 breaths per minute). Zero-padding to 256 points makes the frequency grid denser;
it does not give a 3 s frame the resolving power of a 16 s observation. These are the
implementation settings of the benchmark branch, not a claim that this configuration is
universally optimal. Each modified frame is inverse transformed;
overlapping samples are accumulated and divided by their overlap counts. Because the mask modifies
spectral magnitudes while reusing mixture phase, the modified coefficients need not be exactly
consistent with any single time-domain signal; this is a general limitation of magnitude-only
masking and modified-STFT reconstruction \citep{griffin1984signal,muller2021fundamentals}.

\subsection{NMF, sparsity, and random-body-movement mitigation}

NMF approximates a non-negative magnitude spectrogram $V\in\mathbb{R}_{\geq0}^{F\times T}$ as
\begin{equation}
  V\approx\nmfW\nmfH,
  \qquad \nmfW\in\mathbb{R}_{\geq0}^{F\times K},
  \quad \nmfH\in\mathbb{R}_{\geq0}^{K\times T}.
  \label{eq:ch2_nmf}
\end{equation}
The columns of $\nmfW$ act as spectral patterns and the rows of $\nmfH$ describe their
time-varying activations. With a fixed rank and non-negativity constraints, NMF can express
recurring spectrotemporal structure without a hand-written template for every artifact. It does
not, however, identify physiological sources by mathematical necessity: component interpretation
and removal depend on the data, objective, initialization, regularization, and decision rule.
Optimization is generally non-convex in the joint factors and iterative algorithms may stop at a
maximum-iteration limit, so solver settings and convergence records are part of the method rather
than incidental implementation details.

Sparsity provides one possible description of an activation pattern. For a nonzero vector of
length $T>1$, $h\in\mathbb{R}_{\geq0}^{T}$, Hoyer sparsity can be written
\begin{equation}
  \hoyerSpar(h)=
  \frac{\sqrt{T}-\lVert h\rVert_1/\lVert h\rVert_2}{\sqrt{T}-1}.
  \label{eq:ch2_hoyer}
\end{equation}
It approaches zero for a uniform nonnegative vector and one for a vector concentrated in a small
number of coordinates. A transient, high-amplitude activation may therefore be a candidate motion
component, but sparsity alone is not a reliable artifact label: breathing itself can be irregular,
and an infant's movement can persist across multiple frames. The Beltrão benchmark uses an
NMF-based component rejection strategy, and this thesis retains rank-11 NMF as a comparison branch
while auditing the dependence on cohort, solver, and threshold choices in later chapters. Hoyer
sparsity is one possible descriptor; its discussion here does not imply that every NMF branch uses
that rejection rule or that the local implementation exactly reproduces the published solver
\citep{beltrao2022contactless,hoyer2004non}.

\subsection{Harmonic--percussive separation as a structural alternative}

Harmonic--percussive source separation (HPSS) originated in audio analysis. Its premise is
geometric: sustained pitched sounds tend to form approximately horizontal ridges in a
time--frequency magnitude spectrogram, whereas brief broadband impacts tend to form vertical
streaks. Median filtering along time and frequency supplies robust local estimates of these two
textures. For magnitude $|X[k,m]|$, the estimates are
\begin{align}
  \hpssHarm[k,m] &= \operatorname{median}_{r\in\mathcal{N}_t(m)}|X[k,r]|,\\
  \hpssPerc[k,m] &= \operatorname{median}_{q\in\mathcal{N}_f(k)}|X[q,m]|,
  \label{eq:ch2_hpss_medians}
\end{align}
where $\mathcal{N}_t$ and $\mathcal{N}_f$ are temporal and frequency neighborhoods. A soft
Wiener-style mask gives greater weight to the time-persistent estimate:
\begin{equation}
  \hpssMask[k,m]=
  \frac{\hpssHarm[k,m]^{\hpssGamma}}
  {\hpssHarm[k,m]^{\hpssGamma}+\hpssPerc[k,m]^{\hpssGamma}+\hpssEps}.
  \label{eq:ch2_hpss_mask}
\end{equation}
The complementary mask can retain the percussive texture; the soft form avoids the abrupt decisions
of a binary mask. Median-filter HPSS and its diffusion-based predecessors were designed for audio,
not for a physiological source-separation guarantee \citep{ono2008separation,fitzgerald2010harmonic,muller2021fundamentals}.

In this thesis, the terminology “harmonic” and “percussive” describes spectrogram geometry only.
It does not assert that $\hpssHarm$ is pure respiration or that $\hpssPerc$ is pure body movement.
An infant's breathing can be nonstationary, shallow, or interrupted; gross movement can also be
slow or sustained. HPSS offers a deterministic structural prior with low algorithmic complexity,
whereas NMF learns a factorization from each spectrogram. The comparison is scientifically useful
precisely because these priors differ, not because their components are known biological sources.
The locked v1 mask is applied on a two-sided, frequency-shifted STFT and is denoted $\hpssMask$.
Later band-restricted or directional variants use the distinct symbol $\hpssMaskBand$; their
definitions must not be interchanged.

The broader HPSS literature explores several extensions. Complementary diffusion and anisotropic
smoothness replace fixed medians with directional regularization; margin-based methods introduce
an ambiguous residual rather than forcing every bin into one of two classes. Multi-resolution
methods combine a longer window for narrow, slowly changing ridges with a shorter window for
rapid events. Kernel additive models generalize the permitted time--frequency directions, and
phase-based methods use instantaneous-frequency or group-delay structure beyond magnitude alone.
Online methods update estimates frame by frame to avoid future-frame dependence. These approaches
clarify design options, but each adds assumptions or computational requirements that must be tested
for respiratory data rather than imported as if audio and neonatal radar shared the same source
statistics \citep{ono2008realtime,driedger2014extending,tachibana2014two,tachibana2014harmonic,fitzgerald2014harmonic,cano2014phase,canadas2015online}.

After time--frequency reconstruction, the local benchmark implementation applies a Kaiser-window
FIR bandpass from $0.30$ to $3.00\,\mathrm{Hz}$ (18--180 breaths per minute) with Kaiser
parameter $\beta=6.5$. Its 64-tap length is marked as inferred in the repository and should not
be attributed to a fully specified filter protocol in the source paper. Forward--backward
filtering is applied separately to the real and imaginary parts, so the resulting filter is
zero-phase but offline and noncausal. It is an implementation choice for retrospective rate
estimation, not a real-time filter with zero latency; finite-window filtering also requires
attention to edge effects.

\section{Harmonic Rate Estimation and Temporal Tracking}
\label{sec:ch2_tracking}

\subsection{From a coarse period to a harmonic estimator}

Respiratory rate estimation is a fundamental-frequency problem. The desired fundamental may be
weaker than a harmonic, and a short window may contain only a small number of cycles. Classical
pitch-estimation methods therefore combine local periodicity with evidence at integer multiples of
a candidate fundamental. The harmonic product spectrum and period-histogram approaches are early
examples of using multiple harmonics to resolve ambiguous peaks; the YIN difference-function
family instead emphasizes periodicity in the time domain and reduces octave errors through
normalized difference structure \citep{schroeder1968period,cheveigne2002yin}.

For a filtered radar trace $\chestDispEst$, an $L_k$-harmonic sinusoidal model at candidate
frequency $f$ can be expressed as
\begin{equation}
  \chestDispEst(t_i) \approx
  \sum_{\ell=1}^{\nlsOrder}
  \left[a_\ell\cos(2\pi\ell f t_i)+b_\ell\sin(2\pi\ell f t_i)\right]
  +\epsilon_i, \qquad a_\ell,b_\ell,\epsilon_i\in\mathbb{C}.
  \label{eq:ch2_harmonic_model}
\end{equation}
For each candidate $f$, the sinusoid coefficients can be estimated by least squares. The profile
cost $\nlsCost$ is the sum of squared residual magnitudes after optimizing those coefficients;
the estimate is the frequency that minimizes this cost over the allowed grid. Under suitable
orthogonality and normalization conditions, this objective can be related to a weighted sum of
spectral power at
$f,2f,\ldots,\nlsOrder f$. The name “NLS” in the benchmark branch should not be read as proof that
the implementation numerically minimizes this joint sinusoidal residual at every candidate: the
repository implementation instead scores a local frequency grid by summing interpolated folded-PSD
power at the first two harmonics, centered on a coarse autocorrelation estimate (with a spectral
fallback when that estimate is unavailable). This is a discrete harmonic-score approximation, and
the exact grid and PSD interpolation settings are implementation choices. Coarse autocorrelation
limits the search range but cannot guarantee that a stronger cardiac or harmonic peak is excluded.
The locked baseline uses a two-harmonic order for this stage, while adaptive-order approaches form
a related but distinct line of work \citep{christensen2008multi,beltrao2023adaptive}.

No single estimator resolves every ambiguity. Autocorrelation can select a multiple of the true
period when cycles are irregular or waveforms are asymmetric. A spectral maximum can favor the
second harmonic. Harmonic summation may reward a candidate whose multiples align with unrelated
cardiac or motion energy. Radar-specific methods therefore examine relationships among harmonic
peaks and between cardiopulmonary components, while state-space models use temporal continuity to
discourage isolated jumps. These strategies supply complementary evidence; continuity cannot
correct a persistent wrong candidate if the emission evidence and transition prior jointly favor
it \citep{naishadham2016estimation,le2020multivariate,beltrao2023adaptive}.

\subsection{Candidate tracking with discrete state-space models}

Let $\viterbiState_t$ index a discrete set of candidate respiratory rates at time $t$. A hidden
Markov model assigns an emission score to each candidate and a transition probability from the
previous state. In log form, a maximum a posteriori path solves
\begin{equation}
  \hat q_{1:T}=\arg\max_{q_{1:T}}
  \left[\log p(q_1)+\sum_{t=1}^{T}\log p(y_t\mid q_t)
  +\sum_{t=2}^{T}\log p(q_t\mid q_{t-1})\right],
  \label{eq:ch2_hmm}
\end{equation}
where $y_t$ denotes the spectral or waveform evidence at frame $t$. The Viterbi algorithm computes
the highest-scoring path efficiently; forward filtering instead propagates a probability
distribution using observations available up to the current time. Forward--backward smoothing uses
future observations and is therefore appropriate for offline analysis but not for a strictly
causal online estimate. HMMs and Viterbi decoding are established tools for sequence inference;
applications to pitch, heart rate, and radar motivate candidate tracking when instantaneous peaks
are unreliable \citep{rabiner1989tutorial,mauch2014pyin,temko2017accurate,xie2019unobtrusive,ye2021spectral,arsalan2020improved}.

The state-transition prior encodes a model of plausible temporal change. If it is too weak, the
path can jump between a fundamental and its harmonic. If it is too strong, it can lag genuine
changes or bridge an apnea pause with an apparently plausible rate. The distinction between a
causal filter and an offline smoother is therefore substantive: they have different information
sets and latency properties, even when they share emissions and transition parameters. Harmonic
salience representations and differentiable dynamic programming offer ways to learn compact
emission weights while retaining an explicit state structure
\citep{bittner2017deep,mensch2018differentiable}.

\subsection{Equivariance and the risk of learned shortcuts}

Frequency translation equivariance is a useful design principle for fundamental-frequency
estimation. If all frequency coordinates and corresponding harmonic evidence shift by a constant,
an estimator with this property shifts its output by the same amount instead of preferring an
absolute frequency learned from the training distribution. PESTO develops a transposition-equivariant
objective for pitch estimation, providing a relevant example of building this symmetry into a
learned system \citep{riou2023pesto}. In a radar tracker, such a constraint can reduce reliance on
the marginal distribution of rates in a small training cohort. It does not guarantee correct
estimation when the signal has no respiratory information or when the observation model is outside
the design range.

The general shortcut-learning literature describes how flexible predictors can exploit stable but
non-causal cues that work in a training distribution and fail under a shift. In respiratory-rate
datasets, a model may use monitor smoothing, cohort prevalence, or recording-specific structure in
place of radar evidence. This is a risk to test, not an assumed explanation for every neural model.
The thesis therefore compares compact physically constrained operators with higher-capacity
regressors under patient-held-out evaluation, a second reference, frequency shifts, and an external
hardware test. Those experiments assess whether the observed predictions track the signal across
conditions; the literature alone cannot establish the outcome of those tests
\citep{geirhos2020shortcut,riou2023pesto}.

\section{Signal Quality, Selective Prediction, and Clustered Statistical Evaluation}
\label{sec:ch2_uncertainty_statistics}

\subsection{Signal-quality indices and risk--coverage analysis}

A point estimate does not say whether its supporting signal is trustworthy. Signal-quality indices
(SQIs) summarize properties such as morphology, periodicity, spectral concentration, or agreement
among independent estimates. Their use in ECG, photoplethysmography, and respiratory-rate
estimation provides a foundation for radar quality assessment, but an SQI must be validated against
the failure modes of its own sensor and population. A rule that rejects broadband spectra may
identify sensor noise, yet it may also reject low-amplitude or interrupted breathing. Conversely,
a plausible rate inside a physiological range does not establish that the radar captured breathing
at all \citep{orphanidou2015signal,birrenkott2018robust}.

Selective prediction formalizes the trade-off between how many observations are retained and the
error among those retained. Let $s_i$ be an uncertainty score, with larger values indicating less
confidence, and accept observations below a threshold $q$. Coverage is the accepted fraction; risk
is a chosen error measure conditional on acceptance. Varying $q$ produces a risk--coverage curve,
which can be summarized by area-under-curve measures such as AURC or by adjusted variants such as
E-AURC. These summaries describe ranking quality over a range of operating points; they do not
select a clinically acceptable operating threshold, determine the consequences of abstention, or
establish calibration by themselves. Learned reject heads such as SelectiveNet are one approach;
physiological SQIs and model-derived posterior spread are alternatives with different training and
assumption requirements \citep{elyaniv2010foundations,geifman2017selective,geifman2019selectivenet,traub2024overcoming}.

For neonatal monitoring, rejected windows are not interchangeable missing values. If uncertainty
is high during quiet breathing pauses, a quality gate may selectively remove the very periods that
carry safety relevance. If the score is insensitive to sensor corruption, it can retain plausible
but unsupported estimates. Evaluation should consequently include both aggregate risk--coverage
behavior and stratified checks for physiological pauses, movement, and out-of-distribution
corruption. The thesis treats pause rejection as a safety-oriented analysis condition, not as
evidence that a radar-only algorithm can diagnose apnea.

\subsection{Conformal prediction with repeated observations per patient}

Conformal prediction converts calibration residuals into prediction sets or intervals under an
exchangeability condition. In the simplest split-conformal setting, a model is fixed before a
calibration sample of exchangeable units is scored. A finite-sample quantile of those scores then
sets the interval width. The unit of exchangeability matters. NICU records contain many highly
overlapping windows from each infant and repeated sessions from the same infant; the windows are
not independent patient-level draws. Randomly dividing windows between training and calibration,
or pooling all within-infant residuals as if they were separate patients, can therefore yield
misleading coverage estimates for a new infant.

Hierarchical methods address data with observations nested within groups. Dunn, Wasserman, and
Ramdas study two-layer prediction sets and compare CDF pooling with single or repeated subsampling.
Those methods clarify the choice between use of all within-group observations and stronger
distribution-free guarantees, but a particular implementation must match its own assumptions and
calibration unit \citep{dunn2023distribution}. The two-level patient-quantile procedures evaluated
later in this thesis are related to this literature; they are not interchangeable with every
procedure in Dunn et al. C4/C5 first form a within-patient empirical quantile of calibration-window
errors and then calibrate an outer order statistic across patient-level scores (C5 applies this to
uncertainty-normalized errors). Their conservative finite-sample $(1-\alpha)^2$ target requires the
exchangeability assumptions at both calibration levels; it is not Dunn et al.'s single
$(1-\alpha)$ guarantee. The overlapping windows in this dataset do not by themselves establish
within-patient exchangeability, so this bound is a conditional theoretical result, not a claim of
verified coverage for the present cohort. The empirical head-to-head comparison and its exact
coverage semantics are reported in Chapter 5.

Finite calibration size imposes a further restriction. With $K_{\mathrm{cal}}$ exchangeable
patient-level calibration scores, the smallest nonzero attainable miscoverage is
\begin{equation}
  \alpha_{\min}=\frac{1}{K_{\mathrm{cal}}+1}.
  \label{eq:ch2_finite_k}
\end{equation}
In leave-one-patient-out folds from 11 infants, $K_{\mathrm{cal}}=10$ and
$\alpha_{\min}=1/11\approx0.0909$. At a nominal 95\% target, the split-conformal order statistic
has index $\lceil(10+1)0.95\rceil=11$, beyond the ten available calibration patients, so the
standard finite-sample quantile is infinite. At least 19 calibration patients are needed for a
non-vacuous 95\% quantile under this rank calculation. This is a statement about the finite
calibration design, not a claim that 19 infants alone suffice for clinical validation or reliable
subgroup coverage.

{\sloppy
Other conformal frameworks address departures from exchangeability through bounded coverage-gap
results, risk control, or online adaptation. Their guarantees depend on specified assumptions,
weighting choices, loss bounds, and access to feedback. In particular, an online method that updates
from observed outcomes cannot be operated in the same way when contemporaneous reference labels
are unavailable at the intended bedside. The distinction between formal marginal coverage,
patient-conditional performance, and operational feasibility is central when assessing a
prediction interval for a new infant \citep{barber2023conformal,angelopoulos2024conformal,gibbs2024conformal}.
\par}

\subsection{Agreement and inference with clustered measurements}

Respiratory-rate estimates are produced repeatedly within each session and infant. Treating every
overlapping window as an independent observation overstates the effective sample size and can make
uncertainty intervals too narrow. Pearson correlation is not an agreement measure: two methods
can be highly correlated while differing systematically in level or scale. Bland--Altman analysis
instead examines paired differences and their dispersion. For repeated observations, the limits
of agreement must account for within-person clustering and, where relevant, additional nesting by
session. Repeated-measures extensions and mixed-effects formulations provide tools for estimating
those components and their uncertainty \citep{bland1986statistical,bland1999measuring,bland2007agreement,zou2013confidence,parker2016application}.

The repeated-measures correlation $\rWithin$ estimates a common within-subject association after
separating subject-level differences; this differs from pooling all windows and computing one
ordinary correlation. The unweighted mean of per-session correlations is another distinct
summary and must be named separately. This distinction matters when infants have different numbers
of valid windows or occupy different respiratory-rate ranges \citep{bland1995calculating,bakdash2017repeated}.

Finally, a nonsignificant difference does not establish equivalence. The two one-sided tests
procedure (TOST) evaluates whether the effect lies within a prespecified interval
$[-\tostMargin,+\tostMargin]$; equivalently, the corresponding 90\% confidence interval must
fall inside the margin for the conventional paired-test formulation. The thesis uses a
predefined margin of 0.50 breaths per minute for its NMF--HPSS comparison, so conclusions are
bounded by that tolerance and the evaluated cohorts. Equivalence testing answers a different
question from testing whether the mean difference is zero, and neither result establishes
generalization to a new site or clinical benefit \citep{schuirmann1987comparison,lakens2017equivalence}.

\section{Synthesis and Positioning of the Thesis}
\label{sec:ch2_synthesis}

The source matrix contains six thematic clusters. Their value lies in the connections between
them: the radar physics explains harmonic ambiguity; the separation literature supplies alternative
structural priors; fundamental-frequency methods provide local candidates; sequence models impose
temporal structure; uncertainty methods characterize when predictions should be trusted; and
cluster-aware statistics determine how evidence across windows and infants should be summarized.
The contributions and limits of each cluster are summarized below.

\begin{table}[htbp]
  \centering
  \footnotesize
  \caption{Thematic synthesis of the 71-reference source matrix and its relationship to the thesis.}
  \label{tab:ch2_literature_synthesis}
  \begin{tabularx}{\textwidth}{@{}>{\raggedright\arraybackslash}p{2.55cm}>{\raggedright\arraybackslash}p{2.15cm}>{\raggedright\arraybackslash}X >{\raggedright\arraybackslash}X@{}}
    \toprule
    Theme & Source set & Main methodological contribution & Relevance and boundary in this thesis \\
    \midrule
    Neonatal radar modalities & Cluster 1 ($n=13$) & CW, IR-UWB, FMCW, optical, and multimodal neonatal sensing; clinical feasibility and review evidence. & Motivates contactless sensing and defines the distinction between a clinical cohort, an adult dataset, and a stationary phantom. \\
    CW phase and I/Q & Cluster 2 ($n=12$) & Quadrature geometry, phase demodulation, motion cancellation, receiver calibration, and Bessel harmonic structure. & Establishes the measurement model; no demodulator removes every source of motion or hardware error. \\
    Spectrogram separation & Cluster 3 ($n=12$) & NMF sparsity, median HPSS, diffusion, margins, multi-resolution, phase, and online variants. & Grounds a controlled NMF--HPSS comparison; audio texture labels are not biological source labels. \\
    Fundamental and sequence tracking & Cluster 4 ($n=11$) & Harmonic estimation, autocorrelation, Viterbi, Kalman, and radar-specific tracking. & Explains candidate ambiguity and temporal priors; smoothing assumptions and latency must be reported. \\
    SQI and clustered agreement & Cluster 5 ($n=14$) & Signal-quality assessment, risk--coverage metrics, repeated-measures correlation and agreement, and TOST. & Motivates selective evaluation and patient/session-aware inference rather than treating windows as independent. \\
    Structured learning and conformal inference & Cluster 6 ($n=9$) & Differentiable dynamic programming, equivariant pitch estimation, shortcut learning, and hierarchical or online conformal methods. & Positions compact structured models and finite-$K$ limits; guarantees remain conditional on their assumptions and calibration units. \\
    \bottomrule
  \end{tabularx}
\end{table}

The literature therefore supports a coherent methodological program but does not predetermine its
empirical conclusions. Radar physics makes harmonic confusion plausible; it does not show which
artifact dominates a given recording. NMF and HPSS express different assumptions about
time--frequency structure; their relative value must be tested on the same inputs and evaluation
population. Temporal continuity can stabilize a sequence while also carrying a wrong state forward.
Uncertainty scores can rank errors while failing a safety constraint at a particular operating
point. Conformal procedures can provide finite-sample statements only for the calibration unit and
exchangeability structure they actually use. These distinctions motivate the thesis's sequence of
forensic checks, controlled model comparisons, selective-risk analyses, and zero-touch external
test.

In the next chapter, these foundations are applied to the benchmark reconstruction and its
protocol uncertainties. Chapter 4 then evaluates front-end and temporal-model choices; Chapter 5
examines selective uncertainty and patient-level conformal behavior; Chapter 6 evaluates transfer
to a separate 24.2 GHz thorax simulator. Throughout, exact reproduction of the original Beltrão
cohort and denominator remains blocked by unresolved source details, and stationary phantom
transfer is treated as a physical-domain result rather than evidence of bedside clinical utility.
